\documentclass[10pt,a4paper]{article}
\usepackage{amsmath,amssymb,amsthm,graphicx}
% Minimal retained preamble for S7 SIGMA article cards taken from
% sources/sigma-2025-023-macdonald-identities-weyl-kac-denominator-fo/pinned-original-source.tex
% (doi:10.3842/SIGMA.2025.023, SIGMA 21 (2025), CC BY 4.0).
%
% Omitted and why: the publisher documentclass{sigma} conditional block (the
% sigma class is unavailable in the toolchain; the assembler supplies the
% article class) and the publisher frontmatter commands; no page geometry or
% font setup is carried so it cannot fight the forced xelatex engine. The
% theorem declarations (author capitalized environment names, the family
% sharing one per-section counter anchored on Theorem) and the per-section
% equation numbering are copied verbatim from the pinned source; graphicx is
% carried because a counted span uses raisebox; no author shorthand macro is
% used inside the counted spans, so none is carried.
\usepackage{amsmath,amssymb,amsthm}
\usepackage{graphicx}
\numberwithin{equation}{section}
\newtheorem{Theorem}{Theorem}[section]
\newtheorem{Corollary}[Theorem]{Corollary}
\newtheorem{Lemma}[Theorem]{Lemma}
\newtheorem{Proposition}[Theorem]{Proposition}
 { \theoremstyle{definition}
\newtheorem{Definition}[Theorem]{Definition}
\newtheorem{Note}[Theorem]{Note}
\newtheorem{Example}[Theorem]{Example}
\newtheorem{Remark}[Theorem]{Remark} }
\providecommand{\qed}{\ensuremath{\square}}
\title{Dominance of the shifted weight along minimal-length coset representatives}
\author{Lecouvey, Cédric, Wahiche, David}
\date{doi:10.3842/SIGMA.2025.023; CC BY 4.0}
\begin{document}
\maketitle
\setcounter{section}{4}
\setcounter{Theorem}{3}
\setcounter{equation}{1}
\noindent\textit{Editorial note.} All delivered mathematical text is the author's own, copied from the pinned source file sources/sigma-2025-023-macdonald-identities-weyl-kac-denominator-fo/pinned-original-source.tex; no statement or proof step was rewritten, reordered or removed. The counted span is the author Lemma environment and its own proof environment. The span contains no citation and no reference, so no stub and no bibliography is retained; the author's own qed closing (qedhere inside the final display, tag*\{\textbackslash{}qed\}, \textbackslash{}renewcommand\{\textbackslash{}qed\}\{\}) is kept verbatim.\par\medskip
\noindent\textit{Editorial note.} Printed numbers verified against the source declarations and a full-file walk: the Theorem family counter is per-section (family shared) and the equation counter is per-section (numberwithin); this lemma is the fourth family environment of section 4 (after 4.1 at line 681, 4.2 at 728, 4.3 at 773) and one numbered display (source line 753) precedes the span in section 4, so initial\_counters section 4, Theorem 3, equation 1 reproduce the printed number 4.4. All displays inside the span are unnumbered, so there is no in-span drift.\par\medskip
\begin{Lemma}
	For any dominant weight $\lambda\in M^{\ast}$ and any $u\in W^{\lambda}$
	$($i.e., of minimal length in the cosets of $W/W_{\lambda})$, the weight
	\[
	\eta^{\vee}\lambda+\bigl(u^{-1}(\mathring{\rho})-\mathring{\rho}\bigr)
	\]
	is dominant.
\end{Lemma}

\begin{proof}
	We need to prove that
	\[
	\bigl(\eta^{\vee}\lambda+\bigl(u^{-1}(\mathring{\rho})-\mathring{\rho}\bigr),\alpha_{i}%
	\bigr)=\bigl(\eta^{\vee}(\lambda,\alpha_{i})+\bigl(u^{-1}(\mathring{\rho})-\mathring{\rho
	}\bigr), \alpha_{i}\bigr)\geq0
	\]
	for any $i\in I^{\ast}$. Observe that we have
	\[
	\bigl(\bigl(u^{-1}(\mathring{\rho})-\mathring{\rho}\bigr), \alpha_{i}\bigr) =\bigl(u^{-1}(\mathring{\rho
	}),\alpha_{i}\bigr)-(\mathring{\rho},\alpha_{i})=(\mathring{\rho},u(\alpha
	_{i}))-\frac{1}{2} \Vert \alpha_{i} \Vert ^{2}
	\]
	since \smash{$(\mathring{\rho},\alpha_{i})=\frac{ \Vert \alpha_{i} \Vert
		^{2}}{2}\bigl(\mathring{\rho},\alpha_{i}^{\vee}\bigr)=\frac{ \Vert \alpha
		_{i} \Vert ^{2}}{2}$}. If $u(\alpha_{i})\in R_{+}$, then $(\mathring{\rho
	},u(\alpha_{i}))\geq\frac{1}{2} \Vert \alpha_{i} \Vert ^{2}$ and
	therefore $\bigl(u^{-1}(\mathring{\rho})-\mathring{\rho}\bigr),\alpha_{i})\geq0$. We are
	done because $\lambda$ is dominant and thus $\eta^{\vee}(\lambda,\alpha
	_{i})\geq0$. Now assume $u(\alpha_{i})\in-R_{+}$ is a negative root. Then, we
	know that there exists a reduced expression of $u$ ending by $s_{i}$, that is
	of the form $u=u^{\prime}s_{i}$ with $\ell(u)=\ell(u^{\prime})+1$. Since $u\in
	W^{\lambda}$, this implies that $(\lambda,\alpha_{i})\geq p$ where $p=1$ in
	all affine types except in types \smash{$C_{n}^{(1)}$} and \smash{$G_{2}^{(1)}$} where $p=2$
	and $p=3$, respectively (because $\lambda\in M^{\ast}\cap P_{+}$). Indeed, we
	would have otherwise $u(\lambda)=u^{\prime}s_{i}(\lambda)=u^{\prime}$ and~$u$~would
    not be of minimal of length. We thus have
	\[
	\bigl(\eta^{\vee}\lambda+\bigl(u^{-1}(\mathring{\rho})-\mathring{\rho}\bigr),\alpha_{i}\bigr)\geq
	p\eta^{\vee}-(\mathring{\rho},-u(\alpha_{i}))-\frac{1}{2} \Vert \alpha
	_{i} \Vert ^{2}\geq0,
	\]
	where $\alpha\in R_{+}$. One can then check that for any $\alpha\in R_{+}$,
	we have $(\mathring{\rho},-u(\alpha_{i}))+\frac{1}{2} \Vert \alpha
	_{i} \Vert ^{2}\leq p\eta^{\vee}$ and thus the desired inequality
	\[
	\bigl(\eta^{\vee}\lambda+\bigl(u^{-1}(\mathring{\rho})-\mathring{\rho}\bigr),\alpha_{i}%
	\bigr)\geq0.\qedhere
\tag*{\qed}
\]
\renewcommand{\qed}{}
\end{proof}

\end{document}
